% ---- pg-386 （书页 374） ----

\chapter{因子分析与对结构性协方差矩阵的推断}
\phantomsection\label{chap:9}

\section{引言}
\phantomsection\label{sec:9.1}

因子分析在其全部历史上时时都激起相当激烈的争论. 它的现代起源在 20 世纪早期, K. 皮尔逊(Pearson), C. 斯皮尔曼(Spearman)及其他一些学者, 为定义和测定智力所作的努力. 因为与智力这类概念早有联系, 主要由对心理测量学有兴趣的科学家们, 培育和发展了因子分析. 对几个早期研究的心理学解释的争论以及缺乏强有力的计算工具, 阻碍了它作为统计学方法的最新发展. 高速计算机的出现, 已经重新引发对因子分析的理论和计算方面的兴趣. 原先的大部分技巧已经被摒弃, 而紧随新近的发展, 早期的争论也消退了. 无论如何, 有一点总是对的: 该技术的每一次应用成功与否, 都必须由其自身功绩来检测.

因子分析的实质目的是, 只要可能, 就用几个潜在的但不能观察的随机量去描述许多变量间的协方差关系, 这些随机量叫做\mvbold{因子}. 基本上说, 因子模型的推动是由于下一命题: 假定可用变量间的相关性把它们分组. 这就是说, 假设在一个特定组内的所有变量, 它们之间是高度相关的, 而与不同组中的变量却有相对较小的相关性. 于是可以想象, 各组变量象征有单一的潜在的结构(或因子), 它对观察到的相关负责. 例如, 由斯皮尔曼收集了古典文学、法语、英语、数学以及音乐的一组测验成绩, 从它们的相关性表明存在一个潜在的因子“智力”. 另一组变量, 表示身体健康得分, 只要有效就可以对应为另一个因子. 因子分析企图确认的就是这样一种结构.

可以把因子分析想成是主成分分析的一个扩充. 两者都可看成为在力图逼近协方差矩阵 $\boldsymbol{\Sigma}$. 然而, 基于因子分析模型的近似是更为精细的. 因子分析的主要问题是, 数据是否与一个规定的结构一致.

\section{正交因子模型}
\phantomsection\label{sec:9.2}

有 $p$ 个成分的观测随机向量 $\mathbf{X}$, 有均值 $\boldsymbol{\mu}$ 和协方差矩阵 $\boldsymbol{\Sigma}$. 因子模型要求 $\mathbf{X}$ 是线性依赖于几个不能观测的称之为公共因子的随机变量 $F_1,F_2,\cdots,F_m$ 和 $p$ 个附加的称之为误差或有时也称为\mvbold{特殊因子}的变差源 $\varepsilon_1,\varepsilon_2,\cdots,\varepsilon_p$.\footnote{正如马克斯韦尔(Maxwell, 见文献[12])指出的, 在许多调查中, 倾向认为 $\varepsilon_i$ 是唯一地与各个变量有关的那些测量误差和因子的组合.} 具体的, 因子分析模型是
\[
X_1-\mu_1=l_{11}F_1+l_{12}F_2+\cdots+l_{1m}F_m+\varepsilon_1
\]
